\documentclass[12pt,a4paper]{article}

% ---- Codificacao e idioma ----
\usepackage[utf8]{inputenc}
\usepackage[T1]{fontenc}
\usepackage[brazil]{babel}

% ---- Fonte Times New Roman e espacamento ----
\usepackage{mathptmx} % Times para texto e matematica
\usepackage{setspace}
\onehalfspacing % espacamento 1,5

% ---- Margens ABNT (sup./esq. 3 cm, inf./dir. 2 cm) ----
\usepackage[top=3cm,left=3cm,bottom=2cm,right=2cm]{geometry}

% ---- Pacotes auxiliares ----
\usepackage{amsmath,amssymb}
\usepackage{booktabs}
\usepackage{array}
\usepackage{listings}
\usepackage{xcolor}
\usepackage{graphicx}
\usepackage[hidelinks]{hyperref}
\usepackage{indentfirst}

\lstset{
  basicstyle=\ttfamily\small,
  breaklines=true,
  frame=single,
  columns=fullflexible,
  keepspaces=true,
}

\begin{document}

% =====================================================================
% Capa
% =====================================================================
\begin{titlepage}
\begin{center}
{\large FACULDADE COTEMIG}\\[0.2cm]
{\large Fundamentos Teóricos da Computação}

\vspace{4cm}

{\LARGE\textbf{Implementação de Máquinas Abstratas:}}\\[0.4cm]
{\Large\textbf{AFD, Autômato de Pilha e Máquina de Turing}}

\vspace{3cm}

{\large Trabalho Final da disciplina de Fundamentos Teóricos da
Computação, apresentado ao Professor Júlio César da Silva.}

\vspace{2cm}

{\large\textbf{Integrantes}}\\[0.4cm]
\begin{tabular}{lc}
\toprule
\textbf{Nome} & \textbf{Matrícula} \\
\midrule
Matheus Felipe de Almeida Mota & 72301031 \\
Pedro Simão do Carmo Diniz     & 72301376 \\
Yuri Leal de Sousa             & 72301341 \\
\bottomrule
\end{tabular}

\vfill

{\large Belo Horizonte -- MG\\ 2026}
\end{center}
\end{titlepage}

% =====================================================================
% Resumo
% =====================================================================
\begin{abstract}
\noindent Este relatório descreve a modelagem e a implementação, em C\# (.NET 6 ou superior),
de três máquinas abstratas da teoria da computação: o Autômato Finito
Determinístico (AFD), o Autômato de Pilha (AP) com aceitação por pilha vazia e a
Máquina de Turing (MT). Cada implementação é uma tradução direta da definição
matemática formal. São apresentados os fundamentos teóricos, as decisões de
projeto, os resultados dos testes obrigatórios e uma análise comparativa do poder
computacional das três máquinas.
\end{abstract}
\newpage

% =====================================================================
% Sumário
% =====================================================================
\tableofcontents
\newpage

% =====================================================================
\section{Introdução}
% =====================================================================
A teoria da computação estabelece uma hierarquia de modelos formais capazes de
reconhecer classes de linguagens com diferentes poderes expressivos. Este
trabalho tem por objetivo modelar e implementar três desses modelos --- o
Autômato Finito Determinístico, o Autômato de Pilha e a Máquina de Turing ---
relacionando, em cada caso, a estrutura de dados escolhida com a respectiva
definição matemática.

O código foi escrito em C\# e tem como alvo a plataforma .NET~6 (compatível com
versões superiores via \emph{roll-forward}), organizado em três
projetos independentes (\texttt{Parte1}, \texttt{Parte2} e \texttt{Parte3}), cada
um separado em camadas de \emph{modelo} (a tupla formal e suas estruturas),
\emph{serviço} (simulação e construção das máquinas) e \emph{aplicação}
(entrada e saída). O objetivo pedagógico é evidenciar que cada máquina é apenas a
expressão computacional de sua definição formal, e não uma solução empírica.

% =====================================================================
\section{Fundamentação Teórica}
% =====================================================================

\subsection{Autômato Finito Determinístico (AFD)}
Um AFD é uma 5-tupla $M = (Q, \Sigma, \delta, q_0, F)$, em que $Q$ é um conjunto
finito de estados, $\Sigma$ é o alfabeto de entrada,
$\delta: Q \times \Sigma \to Q$ é a função de transição (total), $q_0 \in Q$ é o
estado inicial e $F \subseteq Q$ é o conjunto de estados de aceitação. Uma cadeia
$w \in \Sigma^*$ é aceita se a computação de $M$ sobre $w$ termina em um estado
$q \in F$.

A linguagem-alvo é $L_1 = \{\, w \in \{a,b\}^* \mid w \text{ termina com } ab \,\}$,
cuja expressão regular equivalente é $(a|b)^*\,ab$. O AFD construído possui:
\[
Q=\{q_0,q_1,q_2\},\quad \Sigma=\{a,b\},\quad q_0=q_0,\quad F=\{q_2\},
\]
com a função de transição dada pela Tabela~\ref{tab:afd}.

\begin{table}[h]
\centering
\caption{Função de transição $\delta$ do AFD para $L_1$.}
\label{tab:afd}
\begin{tabular}{c|cc}
\toprule
$\delta$ & $a$ & $b$ \\
\midrule
$\rightarrow q_0$ & $q_1$ & $q_0$ \\
$q_1$             & $q_1$ & $q_2$ \\
$*\,q_2$          & $q_1$ & $q_0$ \\
\bottomrule
\end{tabular}
\end{table}

\noindent\textbf{Cálculo manual.} Para $w = bab$:
$q_0 \xrightarrow{b} q_0 \xrightarrow{a} q_1 \xrightarrow{b} q_2 \in F$, logo
\textsc{aceita}. Para $w = ba$:
$q_0 \xrightarrow{b} q_0 \xrightarrow{a} q_1 \notin F$, logo \textsc{rejeita}.

\subsection{Autômato de Pilha (AP)}
Um AP é uma 7-tupla $M = (Q, \Sigma, \Gamma, \delta, q_0, Z_0, \varnothing)$, em
que $\Gamma$ é o alfabeto da pilha, $Z_0 \in \Gamma$ é o símbolo inicial da pilha
e $\delta: Q \times (\Sigma \cup \{\varepsilon\}) \times \Gamma \to
\mathcal{P}(Q \times \Gamma^*)$. Adotamos a aceitação \textbf{por pilha vazia}: a
cadeia $w$ é aceita se existe uma computação que consome toda a entrada e termina
com a pilha vazia.

A linguagem-alvo é $L_2 = \{\, a^n b^n \mid n \geq 1 \,\}$. Empilha-se um símbolo
$A$ para cada $a$ lido e desempilha-se um $A$ para cada $b$; ao final, um
$\varepsilon$-movimento remove $Z_0$, esvaziando a pilha.

\noindent\textbf{Cálculo manual 1 (aceitação).} Para $w = aabb$, partindo de $(q_0, aabb, Z)$:
\[
(q_0, aabb, Z) \vdash (q_0, abb, AZ) \vdash (q_0, bb, AAZ) \vdash (q_1, b, AZ)
\vdash (q_1, \varepsilon, Z) \vdash (q_1, \varepsilon, \varepsilon),
\]
pilha vazia com entrada consumida: \textsc{aceita}.

\noindent\textbf{Cálculo manual 2 (rejeição).} Para $w = aab$, partindo de $(q_0, aab, Z)$:
\[
(q_0, aab, Z) \vdash (q_0, ab, AZ) \vdash (q_0, b, AAZ) \vdash (q_1, \varepsilon, AZ).
\]
A entrada foi consumida, mas a pilha ainda contém $AZ$ (sobrou um $A$ sem $b$
correspondente) e não há transição $\delta(q_1, \varepsilon, A)$ definida: a
computação trava sem esvaziar a pilha, logo \textsc{rejeita}.

Como desafio, implementa-se também
$L_3 = \{\, w \in \{a,b\}^* \mid w = w^R,\ |w| \geq 1 \,\}$ (palíndromos), que exige
\emph{não determinismo}: o autômato empilha a primeira metade, ``adivinha'' o
ponto médio (por $\varepsilon$-movimento, no caso par, ou lendo o símbolo central,
no caso ímpar) e compara a segunda metade desempilhando.

\paragraph{Diferença de complexidade entre $L_2$ e $L_3$.}
Embora ambas sejam livres de contexto, as duas linguagens ocupam posições
distintas dentro dessa classe. $L_2$ é uma linguagem livre de contexto
\emph{determinística} (DCFL): em cada configuração há no máximo uma transição
aplicável, pois a própria estrutura da cadeia ($a$'s seguidos de $b$'s) indica,
pelo símbolo lido, se o autômato deve empilhar ou desempilhar; um AP
determinístico (DPDA) a reconhece. Já $L_3$ é o exemplo clássico de LLC
\emph{inerentemente não determinística}: em uma única passada da esquerda para a
direita, o autômato não tem como identificar o ponto médio da cadeia, pois essa
decisão depende de símbolos que ainda não foram lidos. Demonstra-se formalmente
que nenhum DPDA reconhece os palíndromos sobre $\{a,b\}$ \cite{hopcroft}: a
estratégia obrigatória --- empilhar a primeira metade e compará-la com a segunda
--- exige ``adivinhar'' onde a primeira metade termina, poder de escolha que um
autômato determinístico não possui. Portanto $L_2 \in \mathrm{DCFL}$, enquanto
$L_3 \in \mathrm{LLC} \setminus \mathrm{DCFL}$ --- e é exatamente por isso que o
simulador implementado precisa explorar múltiplas computações alternativas
(busca em profundidade com retrocesso) para $L_3$, enquanto para $L_2$ uma única
computação determinística basta.

\subsection{Máquina de Turing (MT)}
Uma MT é uma 7-tupla $M = (Q, \Sigma, \Gamma, \delta, q_0, q_{accept}, q_{reject})$,
com $\Sigma \subseteq \Gamma$, $\sqcup \in \Gamma \setminus \Sigma$ e
$\delta: Q \times \Gamma \to Q \times \Gamma \times \{L, R\}$, sendo
$q_{accept} \neq q_{reject}$.

A linguagem-alvo é $L_4 = \{\, a^n b^n c^n \mid n \geq 1 \,\}$. A estratégia é de
\emph{marcação}: a cada ciclo marca-se um $a$ (para $X$), um $b$ (para $Y$) e um
$c$ (para $Z$), retornando à esquerda e repetindo até que não reste nenhum $a$;
em seguida verifica-se que só restam $Y$ e $Z$ antes do branco.

\noindent\textbf{Cálculo manual 1 (aceitação, esboço) para $w = abc$:}
\[
q_0\,\underline{a}bc \vdash X q_1\,\underline{b}c \vdash XY q_2\,\underline{c}
\vdash X q_3\,\underline{Y}Z \vdash \dots \vdash q_{accept},
\]
todos os símbolos marcados como $X,Y,Z$: \textsc{aceita}.

\noindent\textbf{Cálculo manual 2 (rejeição) para $w = ab$:}
\[
q_0\,\underline{a}b \vdash X q_1\,\underline{b} \vdash XY q_2\,\underline{\sqcup}.
\]
Em $q_2$ a máquina procura um $c$ não marcado, mas encontra o branco $\sqcup$;
como $\delta(q_2, \sqcup)$ não está definida, a computação para em estado
não-aceitador, comportamento equivalente a $q_{reject}$: \textsc{rejeita}.

O desafio implementa uma segunda MT que \textbf{computa} $f(n) = n+1$ em unário:
a entrada possui $n$ símbolos $1$ e, ao final, a fita contém $n+1$ símbolos $1$
(percorre-se a fita à direita até o primeiro branco, que é substituído por $1$).

% =====================================================================
\section{Descrição da Implementação}
% =====================================================================
Cada projeto é um aplicativo de console .NET, organizado em \texttt{Models/},
\texttt{Services/} e \texttt{Program.cs}, seguindo o princípio de
responsabilidade única e evitando classes aninhadas.

\subsection{Parte 1 --- AFD}
A classe \texttt{AutomatoFinitoDeterministico} expõe os cinco componentes da tupla
com nomes fiéis à definição (\texttt{Estados}, \texttt{Alfabeto},
\texttt{Transicao}, \texttt{EstadoInicial}, \texttt{EstadosAceitacao}). A função
$\delta$ é, conforme exigido, um
\texttt{Dictionary<(string estado, char simbolo), string>}. O método
\texttt{Simular} percorre a cadeia símbolo a símbolo, registra o rastro de estados
e retorna um \texttt{ResultadoSimulacao}; \texttt{Aceitar} é um atalho booleano.
O método \texttt{ExibirDiagrama} imprime a tabela de transições $\delta$ com
molduras em ASCII e larguras de coluna ajustadas dinamicamente ao conteúdo,
funcionando para qualquer AFD carregado do JSON. Os resultados das simulações são
apresentados em uma tabela-resumo alinhada --- uma linha por cadeia, com as colunas
\textsc{cadeia}, \textsc{status} e \textsc{rastro} (mais uma coluna \textsc{motivo}
exibida apenas quando há rejeição com causa registrada) --- seguida de uma linha de
totais (número de cadeias, aceitas e rejeitadas); o \textsc{status} é colorido em
verde/vermelho quando a saída vai para o terminal. O desafio é
atendido por \texttt{AfdJsonLoader}, que carrega qualquer AFD a partir de
\texttt{afd.json} (campos \texttt{estados}, \texttt{alfabeto},
\texttt{estadoInicial}, \texttt{estadosAceitacao} e \texttt{transicoes}),
validando o esquema antes da construção: campos obrigatórios presentes,
símbolos de exatamente um caractere, transições que referenciam apenas estados
de $Q$ e símbolos de $\Sigma$, e ausência de transições duplicadas (que
violariam o determinismo de $\delta$).

\subsection{Parte 2 --- Autômato de Pilha}
A pilha é uma \texttt{Stack<char>} manipulada manualmente. A função $\delta$ é um
\texttt{Dictionary} de \texttt{(string estado, char entrada, char topo)} para uma
lista de transições, permitindo não determinismo; $\varepsilon$-movimentos são
representados por \texttt{'\textbackslash 0'}. O \texttt{SimuladorPilha} explora as
computações por busca em profundidade com conjunto de configurações visitadas
(para evitar laços) e limite de passos. A cada passo é registrada a configuração
instantânea (estado, entrada restante e conteúdo da pilha), exibida em uma
tabela com molduras (colunas \textsc{passo}, \textsc{estado}, \textsc{entrada},
\textsc{pilha} e \textsc{transição}) tanto para a computação aceitadora quanto
para o caminho mais profundo explorado em caso de rejeição; o status
(\textsc{aceita}/\textsc{rejeita}, colorido no terminal) é exibido junto à cadeia
e cada linguagem é encerrada com uma linha de totais.

\subsection{Parte 3 --- Máquina de Turing}
A fita é uma estrutura dinâmica \texttt{Dictionary<int,char>}, com o branco
representado por \texttt{'\_'}. A função $\delta$ é um
\texttt{Dictionary<(string estado, char simbolo), (string novoEstado, char novoSimbolo, char direcao)>}.
O \texttt{SimuladorTuring} mantém um contador e um limite configurável de passos
(evitando laços infinitos), que pode ser definido por argumento de linha de
comando (ex.: \texttt{dotnet run --project Parte3 -- 5000}), e registra, a cada
passo, o estado atual, a fita completa com o símbolo sob o cabeçote entre
colchetes e a posição do cabeçote. Esse histórico é apresentado pela camada de
aplicação em uma tabela com molduras (colunas \textsc{passo}, \textsc{estado},
\textsc{cabeçote} e \textsc{fita}), seguida do status (\textsc{aceita},
\textsc{rejeita} ou \textsc{limite}, colorido no terminal) com o número de passos
e, para a MT que computa $f(n)=n+1$, a fita final; cada seção termina com uma
linha de totais. Há estados explícitos \texttt{qaccept} e \texttt{qreject};
transições indefinidas são tratadas como rejeição.

% =====================================================================
\section{Resultados e Testes}
% =====================================================================

\begin{table}[h]
\centering
\caption{Resultados da Parte 1 (AFD para $L_1$).}
\begin{tabular}{lll}
\toprule
Cadeia & Resultado & Observação \\
\midrule
ab      & ACEITA  & caso mínimo \\
aab     & ACEITA  & termina com ``ab'' (ver nota) \\
bab     & ACEITA  & prefixo irrelevante \\
ababab  & ACEITA  & múltiplas ocorrências \\
ba      & REJEITA & termina com ``a'' \\
$\varepsilon$ & REJEITA & cadeia vazia \\
b       & REJEITA & um símbolo \\
\bottomrule
\end{tabular}
\end{table}

\noindent\textbf{Nota sobre \texttt{aab}.} O enunciado sugere \textsc{rejeita}
para \texttt{aab}, mas a cadeia $aab$ \emph{termina} com o sufixo $ab$ e, portanto,
pertence a $L_1$. A implementação segue rigorosamente a definição formal de $L_1$
(coerente com \texttt{bab}, também aceita), resultando em \textsc{aceita}.

\begin{table}[h]
\centering
\caption{Resultados da Parte 2 --- $L_2 = \{a^n b^n\}$ e $L_3$ (palíndromos).}
\begin{tabular}{llp{4cm}}
\toprule
Cadeia & Resultado & Linguagem \\
\midrule
ab     & ACEITA  & $L_2$, $n=1$ \\
aabb   & ACEITA  & $L_2$, $n=2$ \\
aaabbb & ACEITA  & $L_2$, $n=3$ \\
aab    & REJEITA & $L_2$ \\
abb    & REJEITA & $L_2$ \\
ba     & REJEITA & $L_2$ \\
$\varepsilon$ & REJEITA & $L_2$, $n \geq 1$ \\
abab   & REJEITA & $L_2$ \\
\midrule
a      & ACEITA  & $L_3$ \\
aba    & ACEITA  & $L_3$ \\
abba   & ACEITA  & $L_3$ \\
ab     & REJEITA & $L_3$ \\
aab    & REJEITA & $L_3$ \\
\bottomrule
\end{tabular}
\end{table}

\begin{table}[h]
\centering
\caption{Resultados da Parte 3 --- $L_4 = \{a^n b^n c^n\}$ e $f(n)=n+1$.}
\begin{tabular}{lll}
\toprule
Entrada & Resultado / Saída & Observação \\
\midrule
abc       & ACEITA  & $n=1$ \\
aabbcc    & ACEITA  & $n=2$ \\
aaabbbccc & ACEITA  & $n=3$ \\
aabbc     & REJEITA & quantidades distintas \\
ab        & REJEITA & sem $c$ \\
abcabc    & REJEITA & fora de ordem \\
$\varepsilon$ & REJEITA & $n \geq 1$ \\
\midrule
1     & 11     & $f(1)=2$ \\
111   & 1111   & $f(3)=4$ \\
11111 & 111111 & $f(5)=6$ \\
\bottomrule
\end{tabular}
\end{table}

Todos os casos obrigatórios foram reproduzidos pela execução dos projetos
(\texttt{dotnet run --project ParteX}), com saída completa dos rastros, das
configurações instantâneas e dos passos da fita.

% =====================================================================
\section{Análise Comparativa}
% =====================================================================
As três máquinas formam a base da hierarquia de Chomsky. O AFD reconhece
exatamente as \textbf{linguagens regulares}; sua ``memória'' é finita (apenas o
estado corrente), o que o impede de contar uma quantidade arbitrária de símbolos.
O AP acrescenta uma \textbf{pilha}, reconhecendo as \textbf{linguagens livres de
contexto}, como $L_2 = a^n b^n$ e os palíndromos $L_3$ --- impossíveis para um
AFD. Entretanto, a pilha (acesso LIFO, único ponto de acesso) ainda é insuficiente
para $L_4 = a^n b^n c^n$, que exige duas comparações de contagem independentes. A
MT, com fita infinita de leitura e escrita bidirecional, reconhece as
\textbf{linguagens recursivamente enumeráveis} e ainda \emph{computa funções},
sendo estritamente mais poderosa que os modelos anteriores. Resumidamente:
\[
\text{Regulares} \subsetneq \text{Livres de Contexto} \subsetneq
\text{Recursivamente Enumeráveis}.
\]

% =====================================================================
\section{Questões Reflexivas}
% =====================================================================

\subsection{Parte 1 --- AFD}
\textbf{1. Por que $L_1$ é regular? Expressão regular equivalente.}\\
$L_1$ é regular porque pode ser descrita por uma expressão regular e reconhecida
por um autômato finito com memória limitada a três estados. Basta acompanhar o
sufixo já lido: ``nada relevante'', ``acabei de ler um $a$'' e ``terminei em
$ab$''. Como o reconhecimento depende apenas dos dois últimos símbolos, uma
quantidade finita de informação é suficiente, o que caracteriza precisamente as
linguagens regulares. A expressão regular equivalente é $(a|b)^*ab$: qualquer
prefixo sobre $\{a,b\}$ seguido obrigatoriamente do sufixo $ab$. O teorema de
Kleene garante a equivalência entre expressões regulares e autômatos finitos,
confirmando a regularidade de $L_1$.

\noindent\textbf{2. Descrição formal da 5-tupla e justificativa dos estados.}\\
$M = (\{q_0,q_1,q_2\}, \{a,b\}, \delta, q_0, \{q_2\})$. O estado $q_0$ representa
``nenhum progresso em direção ao sufixo $ab$''; $q_1$ representa ``o último símbolo
lido foi $a$'', ou seja, um candidato a iniciar $ab$; $q_2$ representa ``a cadeia
lida até aqui termina em $ab$'' e é o único estado de aceitação. As transições
foram escolhidas para manter essa invariante: ao ler $b$ em $q_1$ chega-se a $q_2$
(sufixo $ab$ completo); ao ler $a$ retorna-se a $q_1$; e ao ler $b$ fora de $q_1$
volta-se a $q_0$. Essa correspondência direta entre significado e estado torna a
construção verificável.

\noindent\textbf{3. E se $\delta$ não fosse total? Tratamento de símbolos inválidos.}\\
Se $\delta$ não fosse total, existiriam pares (estado, símbolo) sem transição
definida e a computação ficaria ``presa'', sem destino. Em um AFD formal, exige-se
a totalidade, geralmente introduzindo um estado de erro (``estado morto'')
não-aceitador para o qual vão todas as transições ausentes. Na implementação, a
ausência de uma transição é tratada explicitamente: ao encontrar um símbolo fora
de $\Sigma$ ou um par sem destino, a simulação interrompe e classifica a cadeia
como \textsc{rejeita}, registrando o motivo (símbolo fora do alfabeto ou transição
indefinida). Assim, entradas inválidas não causam exceções e o comportamento é
equivalente ao de um estado morto.

\subsection{Parte 2 --- Autômato de Pilha}
\textbf{1. Por que $L_2$ não é reconhecível por AFD? (Lema do Bombeamento)}\\
Suponha, por absurdo, que $L_2 = \{a^n b^n\}$ seja regular. Pelo Lema do
Bombeamento para linguagens regulares, existe um comprimento de bombeamento $p$
tal que toda cadeia $w \in L_2$ com $|w| \geq p$ pode ser escrita como $w = xyz$,
com $|xy| \leq p$, $|y| \geq 1$ e $xy^iz \in L_2$ para todo $i \geq 0$. Tome
$w = a^p b^p$. Como $|xy| \leq p$, o trecho $y$ é composto apenas de $a$'s, digamos
$y = a^k$ com $k \geq 1$. Bombeando com $i = 2$, obtemos $a^{p+k} b^p$, que possui
mais $a$'s do que $b$'s e, portanto, não pertence a $L_2$ --- contradição. Logo
$L_2$ não é regular e não pode ser reconhecida por um AFD, pois exige contar uma
quantidade ilimitada de $a$'s, o que ultrapassa a memória finita de estados.

\noindent\textbf{2. Aceitação por pilha vazia e por estado final são equivalentes?}\\
Sim, são equivalentes em poder de reconhecimento: ambas reconhecem exatamente a
classe das linguagens livres de contexto. A demonstração é construtiva. Dado um AP
que aceita por estado final, constrói-se outro que aceita por pilha vazia
adicionando um novo símbolo de fundo $X_0$ abaixo de $Z_0$ e um estado de
``esvaziamento'': sempre que o autômato original atinge um estado final, o novo
autômato passa a desempilhar tudo por $\varepsilon$-movimentos até esvaziar a
pilha. Reciprocamente, dado um AP que aceita por pilha vazia, adiciona-se um
estado final que só é alcançado, via $\varepsilon$-movimento, quando o símbolo de
fundo especial fica exposto (indicando pilha logicamente vazia). Como cada modelo
pode simular o outro, eles reconhecem a mesma classe de linguagens.

\noindent\textbf{3. Papel do símbolo $Z_0$.}\\
$Z_0$ (representado por \texttt{Z} no código) é o símbolo inicial da pilha,
presente antes da leitura de qualquer símbolo. Ele cumpre dois papéis: serve como
marcador de fundo (``a pilha está logicamente vazia se só resta $Z_0$'') e permite
detectar o momento de finalizar a computação. Em $L_2$, os símbolos $A$ são
empilhados sobre $Z_0$ e removidos pelos $b$'s; quando todos os $A$'s foram
retirados, um $\varepsilon$-movimento $\delta(q_1, \varepsilon, Z) = (q_1,
\varepsilon)$ remove $Z_0$ e esvazia a pilha, configurando a aceitação. Sem $Z_0$
não haveria como distinguir ``pilha vazia por sucesso'' de ``pilha vazia por falta
de símbolo a comparar'', nem como acionar com segurança o passo final.

\subsection{Parte 3 --- Máquina de Turing}
\textbf{1. Por que $L_4$ não é reconhecível por AP? Que propriedade da MT permite?}\\
$L_4 = \{a^n b^n c^n\}$ não é livre de contexto, como mostra o Lema do Bombeamento
para LLCs: para $w = a^p b^p c^p$ e qualquer decomposição $w = uvxyz$ com
$|vxy| \leq p$ e $|vy| \geq 1$, o trecho $vxy$ não consegue cobrir os três blocos
simultaneamente; ao bombear, ao menos um dos símbolos ($a$, $b$ ou $c$) fica em
quantidade diferente dos demais, saindo de $L_4$. A pilha do AP só permite uma
contagem ``casada'' (LIFO): ao comparar $a$'s com $b$'s, ela é consumida e não
resta informação para comparar com os $c$'s. A MT supera essa limitação por
possuir uma fita de leitura e escrita com movimento bidirectional: ela pode
\emph{marcar} símbolos e percorrer a fita repetidas vezes, mantendo duas
contagens independentes --- propriedade essencial para reconhecer $L_4$.

\noindent\textbf{2. Quantos estados a MT de $L_4$ usou? Estratégia de marcação.}\\
A MT de $L_4$ utiliza sete estados: $q_0$ (procura o próximo $a$ não marcado),
$q_1$ (avança até o primeiro $b$ não marcado), $q_2$ (avança até o primeiro $c$
não marcado), $q_3$ (retorna à esquerda até o último $X$), $q_4$ (verificação
final de que só restam $Y$ e $Z$), além de $q_{accept}$ e $q_{reject}$. A
estratégia de marcação é: a cada ciclo, troca-se um $a$ por $X$, o primeiro $b$ por
$Y$ e o primeiro $c$ por $Z$, garantindo que as quantidades sejam consumidas na
mesma proporção; depois retorna-se ao início e repete-se. O processo termina quando
não há mais $a$: então $q_4$ confirma que restam apenas $Y$'s e $Z$'s antes do
branco, aceitando. Qualquer desbalanceamento (um $b$ ou $c$ a mais/menos, ou
símbolo fora de ordem) leva a uma transição indefinida e, portanto, à rejeição.

\noindent\textbf{3. Um computador moderno é mais poderoso que uma MT? (Church-Turing)}\\
Não. Pela Tese de Church-Turing, qualquer função efetivamente calculável pode ser
computada por uma Máquina de Turing; nenhum dispositivo físico realizável computa
algo além do que uma MT computa. Um computador moderno é, na verdade,
\emph{menos} poderoso em um aspecto: possui memória finita, enquanto a MT dispõe de
fita ilimitada. Assim, o computador real é, no máximo, equivalente a uma MT de
fita limitada (um autômato finito muito grande). A vantagem prática dos
computadores está na velocidade e na conveniência, não no poder computacional:
problemas indecidíveis para a MT, como o Problema da Parada, permanecem
indecidíveis para qualquer computador real.

\newpage
% =====================================================================
\section{Conclusão}
% =====================================================================
A implementação das três máquinas reforçou, na prática, a correspondência entre
definição matemática e estrutura de dados: o AFD como um dicionário de transições,
o AP como uma pilha explícita com não determinismo e a MT como uma fita dinâmica
com cabeçote. Observou-se concretamente a hierarquia de poder computacional ---
cada modelo reconhece linguagens que o anterior não alcança --- e compreendeu-se
o papel de elementos como o estado morto, o símbolo de fundo da pilha e a
estratégia de marcação na fita. O exercício de traduzir formalismos em código
tornou tangíveis conceitos centrais da teoria da computação.

\newpage
% =====================================================================
\section*{Referências}
\addcontentsline{toc}{section}{Referências}
% =====================================================================
\begin{thebibliography}{9}
\bibitem{sipser} SIPSER, Michael. \textit{Introdução à Teoria da Computação}.
2.~ed. São Paulo: Cengage Learning, 2007.
\bibitem{hopcroft} HOPCROFT, John E.; MOTWANI, Rajeev; ULLMAN, Jeffrey D.
\textit{Introdução à Teoria de Autômatos, Linguagens e Computação}. 3.~ed. Rio de
Janeiro: Campus/Elsevier, 2003.
\bibitem{menezes} MENEZES, Paulo Blauth. \textit{Linguagens Formais e Autômatos}.
6.~ed. Porto Alegre: Bookman, 2011.
\bibitem{lewis} LEWIS, Harry R.; PAPADIMITRIOU, Christos H.
\textit{Elementos de Teoria da Computação}. 2.~ed. Porto Alegre: Bookman, 2000.
\end{thebibliography}

\end{document}
